% =====================================================================
%  Guia — exemplo de uso do template guide.cls
%  Compilar com XeLaTeX ou LuaLaTeX:
%      xelatex main.tex   (rode 2x para o sumário)
%
%  Para versão em inglês, use:  \documentclass[english]{guide}
% =====================================================================
\documentclass{guide}

% --- Configuração do documento --------------------------------------
\setguidetitle{Guia: Acessar o Console da Huawei Cloud}
\setheadertitle{Huawei Cloud -- Guia de acesso ao Console}
\setdocversion{1.0.0}
\setdocdate{2026-08-05}
\setcovertext{Huawei Technologies CO., LTD}

\begin{document}

% ---- CAPA ----------------------------------------------------------
\makecover

% ---- SUMÁRIO -------------------------------------------------------
\sumario

% ---- CORPO (reinicia a numeração de páginas) -----------------------
\startbody

% =====================================================================
%  Capítulo 1 — Acessar o Console da Huawei Cloud
% =====================================================================
\section{Acessar o Console da Huawei Cloud}

\begin{objetivos}
  \objgeral{Este laboratório apresenta o acesso ao Console da Huawei Cloud,
    desde o login no portal até a navegação básica pelo painel de serviços.
    Ao concluir, você será capaz de autenticar-se na plataforma, selecionar a
    região desejada e localizar os principais serviços do catálogo.}

  \prerequisitos
  \begin{itemize}
    \item Conta IAM ativa na Huawei Cloud (com usuário e senha válidos).
    \item Navegador web atualizado (Google Chrome, Mozilla Firefox ou
          Microsoft Edge).
    \item Acesso à internet e permissão para acessar o endpoint
          \code{console.huaweicloud.com}.
  \end{itemize}
\end{objetivos}

\begin{infobox}
Este guia cobre o fluxo básico de acesso ao Console. Para operações
avançadas, consulte a documentação oficial da Huawei Cloud.
\end{infobox}

\subsection{Acessar e navegar no Console}

\objpratica{Realizar o login no portal da Huawei Cloud e percorrer os
  elementos principais do Console: catálogo de serviços, seletor de região
  e área de notificações.}

\passoapasso
\begin{enumerate}
  \item Acesse a \weblink{https://www.huaweicloud.com/intl/en-us/}{Huawei Cloud}
        e, no canto superior direito, clique em \textbf{Console}.

  \imagemc{../../assets/exemplo-login.png}{Figura 1 -- Tela de login do Console.}

  \item Informe o nome de usuário e a senha da sua conta IAM e clique em
        \textbf{Log In}. Se a conta tiver autenticação em dois fatores habilitada,
        informe o código enviado ao seu dispositivo.

  \item Após o login, o painel principal do Console é exibido. No canto
        superior esquerdo, abra o catálogo de serviços clicando no ícone de
        menu.

  \imagem{../../assets/exemplo-menu.png}

  \item No canto superior direito, selecione a região em que seus recursos
        serão criados (por exemplo, \code{sa-brazil-1}).

  \nota{A região selecionada no Console define onde os novos recursos são
    provisionados. Recursos já existentes em outras regiões continuam visíveis
    apenas quando você alterna para a região correspondente.}

  \begin{aviso}
  A região não pode ser alterada para recursos já provisionados.
  Planeje a região antes de criar seus recursos.
  \end{aviso}
\end{enumerate}

\subsection{Explorar o Console}

\subsubsection{Visão geral do painel}

O painel do Console é organizado em três áreas principais: o catálogo de
serviços (à esquerda), a área de trabalho central e a barra superior, onde
ficam o seletor de região, o seletor de projeto e as notificações da conta.
Use o catálogo para abrir qualquer serviço, como \textbf{Elastic Cloud Server
(ECS)} ou \textbf{Object Storage Service (OBS)}.

\paragraph{Regiões e projetos}
Cada região pode conter múltiplos projetos IAM. O seletor de projeto, ao lado
do seletor de região, permite isolar recursos por equipe ou finalidade. O
parâmetro \param{region} do provedor Terraform deve corresponder exatamente ao
identificador da região escolhida aqui, por exemplo \code{sa-brazil-1}.

\subsubsection{Regiões disponíveis}

A tabela abaixo lista as regiões suportadas para o Brasil e América Latina:

\begin{table}[h]
  \centering
  \small
  \begin{tabular}{lll}
    \toprule
    \textbf{Região} & \textbf{Localidade} & \textbf{Código} \\
    \midrule
    São Paulo & Brasil & \code{sa-brazil-1} \\
    Santiago  & Chile & \code{la-south-2} \\
    \bottomrule
  \end{tabular}
  \caption{Regiões disponíveis na Huawei Cloud.}
\end{table}

\imagemplaceholder{assets/console-regions.png}{Captura de tela do seletor de região no Console}

% =====================================================================
%  Capítulo 2 — Configurar credenciais para a CLI
% =====================================================================
\section{Configurar credenciais para a CLI}

\begin{objetivos}
  \objgeral{Este capítulo mostra como obter um par de chaves de acesso
    (AK/SK) no Console e declarar o provedor da Huawei Cloud em um arquivo
    de configuração, preparando o ambiente para uso com a CLI e com
    ferramentas de Infrastructure as Code.}

  \prerequisitos
  \begin{itemize}
    \item Acesso ao Console da Huawei Cloud concluído no Capítulo 1.
    \item Permissão IAM para criar credenciais de acesso.
    \item \code{terraform} instalado localmente (versão 1.0 ou superior).
  \end{itemize}
\end{objetivos}

\subsection{Criar uma AK/SK}

\objpratica{Gerar um par de Access Key ID / Secret Access Key para autenticar
  ferramentas externas na Huawei Cloud.}

No Console, clique no nome do usuário, no canto superior direito, e abra
\textbf{My Credentials}. Na seção \textbf{Access Keys}, clique em
\textbf{Create Access Key} e baixe o arquivo gerado. Guarde o arquivo em local
seguro: a Secret Access Key é exibida apenas uma vez.

\subsection{Declarar o provider}

\passoapasso
\begin{enumerate}
  \item Crie um arquivo chamado \param{provider.tf} no diretório do seu
        projeto com o seguinte conteúdo:

\begin{codigo}
terraform {
  required_providers {
    huaweicloud = {
      source  = "huaweicloud/huaweicloud"
      version = ">= 1.60.0"
    }
  }
}

provider "huaweicloud" {
  region     = "sa-brazil-1"
  access_key = "SUA_ACCESS_KEY_AQUI"
  secret_key = "SUA_SECRET_KEY_AQUI"
}
\end{codigo}

  Substitua os valores de \param{access\_key} e \param{secret\_key} pelas
  credenciais obtidas na seção anterior.

  \item No terminal, na pasta do projeto, inicialize o diretório de trabalho
        do \code{terraform}:

\begin{codigo}[bash]
terraform init
terraform plan
\end{codigo}

  \nota{Nunca versione o arquivo \param{provider.tf} com credenciais reais.
    Use variáveis de ambiente ou um arquivo de estado ignorado pelo controle
    de versão.}
\end{enumerate}

\subsection{Próximos passos}

Com o provider declarado, você pode começar a definir recursos no
\code{terraform}. Para conhecer todos os recursos disponíveis, consulte a
documentação oficial em
\weblink{https://registry.terraform.io/providers/huaweicloud/huaweicloud/latest/docs}{Huawei Cloud Provider}.

\begin{dica}
Use \code{terraform validate} para verificar a sintaxe antes de aplicar
as alterações.
\end{dica}

\badge{Exemplo}

\nota{Este guia é um exemplo de uso do template \code{guide.cls};
  ajuste os textos, imagens e blocos de código conforme o fluxo real do
  seu guia.}

\section{Histórico de versões}

\begin{changelog}
  \changelogentry{1.0.0}{2026-08-05}{
    \item Versão inicial.
    \item Adicionados callout boxes e badge.
  }
\end{changelog}

\end{document}
